\section{A condition that checks every group of students}\label{ch09:sec:condition}
We use the model of Chapter~\ref{ch:matching}: a finite bipartite graph (Definition~\ref{ch08:def:bipartite}) with the set $L$ of students and the set $R$ of tasks as its parts, and an edge $s$--$t$ when student $s$ accepts task $t$. A matching (Definition~\ref{ch08:def:matching}) is a set of edges no two of which share an endpoint, and it saturates $L$ when every student lies on one of its edges; such a matching is exactly a full assignment. For a set $S\subseteq L$, $N(S)$ is the set of tasks acceptable to at least one student in $S$ (Section~\ref{ch08:sec:shortage}; $N$ stands for \emph{neighbors}), and $N(x)$ means $N(\{x\})$.

Section~\ref{ch08:sec:shortage} showed that a \emph{shortage set}, a group $S$ of students with $|N(S)|<|S|$, makes a full assignment impossible. Hall's theorem says that this is the only thing that can go wrong. It concerns the condition
\[
|N(S)|\ge |S|\qquad\text{for every subset }S\subseteq L.
\tag{H}\label{eq:hall}
\]
In words: every group of students accepts, between them, at least as many tasks as there are students in the group. We call \eqref{eq:hall} \emph{Hall's condition}; it is the group condition in the solution to Exercise~\ref{ex:2.5}.

Why must we check every group? Checking each student on their own misses students who compete for the same few tasks, and checking only the whole class misses a smaller group that is short of tasks. For example, let
\[
N(a)=N(b)=\{1\},\qquad N(c)=\{2,3\}.
\]
Every student accepts at least one task, and the whole class $L=\{a,b,c\}$ accepts three tasks, as many as there are students. Yet the group $\{a,b\}$ accepts only task~$1$, and two students cannot both receive it, so no full assignment exists.

So read \eqref{eq:hall} as a list of tests, one for each group $S$: collect into $N(S)$ every task that some member of $S$ accepts (a task accepted by several members counts once, since it can be given to only one of them), and compare $|N(S)|$ with $|S|$. With $n$ students, a group is formed by deciding for each student separately whether they belong to it, so there are $2\cdot2\cdots2=2^n$ groups, counting the empty group and the whole class; for three students there are eight tests. The empty group has $N(\varnothing)=\varnothing$, so its test $0\ge0$ always holds, and a single student $x$ tests $|N(x)|\ge1$.

\par\noindent\begin{minipage}{\linewidth}
\begin{theorem}[Hall's matching theorem]\label{thm:hall}
Let $G$ be a finite bipartite graph with parts $L$ and $R$. Then $G$ has a matching that saturates $L$ if and only if $|N(S)|\ge|S|$ for every subset $S\subseteq L$.
\end{theorem}
\end{minipage}\par

The theorem makes two claims. If a full assignment exists, then \eqref{eq:hall} holds: in the language of Section~\ref{ch02:sec:implication}, Hall's condition is \emph{necessary}. If \eqref{eq:hall} holds, then a full assignment exists: the condition is also \emph{sufficient}.

We prove the sufficient direction twice: once with the search procedure of Chapter~\ref{ch:matching}, which actually constructs an assignment, and once by strong induction (Chapter~\ref{ch:induction}). Philip Hall published the theorem in 1935, and it has become one of the basic tools of combinatorics, the study of finite arrangements and how to count them. For further reading on graphs and matchings, see~\cite{mcs}.
